\documentclass[12pt,a4paper]{article}
\usepackage[utf8]{vietnam}
\usepackage{amsmath,amssymb}
\usepackage{geometry}
\geometry{
  top=12mm,
  bottom=12mm,
  left=14mm,
  right=14mm,
  footskip=6mm,
  headheight=14pt
}

\usepackage{tikz}
\usetikzlibrary{arrows.meta,calc,shapes.geometric}
\usepackage[most]{tcolorbox}
\usepackage{fontawesome}
\usepackage{fancyhdr}
\usepackage{enumitem}

\pagestyle{fancy}
\fancyhf{}
\fancyhead[L]{\footnotesize\color{teal!80!black}\faBook\ \textbf{SỔ TAY KIẾN THỨC TOÁN 3}}
\fancyhead[R]{\footnotesize\color{teal!80!black}\textbf{Lớp Toán Thầy Lĩnh -- loptoanthaylinh.vn}}
\fancyfoot[C]{\footnotesize\color{gray}-- Sổ tay kiến thức: Bài 06. Gấp một số lên một số lần --}
\renewcommand{\headrulewidth}{0.4pt}
\renewcommand{\footrulewidth}{0pt}

\setlength{\parindent}{0pt}
\setlength{\parskip}{0pt}

\begin{document}

% --- BANNER BÀI HỌC ---
\begin{tcolorbox}[
  enhanced,
  colback=teal!8!white,
  colframe=teal!70!black,
  arc=3mm, boxrule=1.2pt,
  left=3mm, right=3mm, top=2mm, bottom=2mm,
  halign=center
]
  {\LARGE\bfseries\color{teal!90!black}\faBookmark\ BÀI 06: GẤP MỘT SỐ LÊN MỘT SỐ LẦN}\\[1mm]
  {\itshape\footnotesize\color{black!75} Quy tắc nhân số đó với số lần, phân biệt bản chất 'gấp lên' và 'nhiều hơn' trên 1 trang A4}
\end{tcolorbox}

\vspace{-2mm}

\noindent
\begin{minipage}[t]{0.49\textwidth}\vspace{0pt}
\begin{tcolorbox}[
  enhanced,
  colback=blue!4!white,
  colframe=blue!75!black,
  arc=2mm, boxrule=0.8pt,
  left=2mm, right=2mm, top=1.8mm, bottom=1.8mm,
  title={\footnotesize\bfseries\color{white}\faMagic\ 1. Quy tắc Gấp lên một số lần}
]
{\footnotesize
\begin{itemize}\setlength{\itemsep}{1pt}\setlength{\leftskip}{-2.5mm}
  \item \textbf{Quy tắc vàng:} Muốn gấp một số lên một số lần, ta lấy \textbf{số đó nhân với số lần}:
  $$\text{Số mới} = \text{Số đã cho} \times \text{Số lần}$$
  \item \textit{Ví dụ:} Đoạn thẳng thứ nhất dài $3\text{ cm}$, đoạn thẳng thứ hai gấp $4$ lần đoạn thứ nhất. Đoạn thứ hai dài: $3 \times 4 = 12\text{ cm}$.
  \item Đoạn thẳng mới gồm các đoạn bằng nhau ghép lại.
\end{itemize}
}
\end{tcolorbox}
\end{minipage}
\hfill
\begin{minipage}[t]{0.49\textwidth}\vspace{0pt}
\begin{tcolorbox}[
  enhanced,
  colback=teal!4!white,
  colframe=teal!75!black,
  arc=2mm, boxrule=0.8pt,
  left=2mm, right=2mm, top=1.8mm, bottom=1.8mm,
  title={\footnotesize\bfseries\color{white}\faBalanceScale\ 2. Phân biệt 'Gấp' \& 'Nhiều hơn'}
]
{\footnotesize
\begin{itemize}\setlength{\itemsep}{1pt}\setlength{\leftskip}{-2.5mm}
  \item \textbf{GẤP $n$ lần} $\implies$ Thực hiện phép \textbf{NHÂN} $(\times)$:
  $$a \text{ gấp } n \text{ lần} \implies a \times n$$
  \item \textbf{NHIỀU HƠN $n$ đơn vị} $\implies$ Thực hiện phép \textbf{CỘNG} $(+)$:
  $$a \text{ nhiều hơn } n \text{ đơn vị} \implies a + n$$
  \item \textit{Ví dụ đối chiếu với số 6 và 3:}
  \begin{itemize}\setlength{\itemsep}{0.5pt}\setlength{\leftskip}{-2.5mm}
    \item $6$ gấp $3$ lần: $6 \times 3 = \mathbf{18}$.
    \item $6$ nhiều hơn $3$ đơn vị: $6 + 3 = \mathbf{9}$.
  \end{itemize}
\end{itemize}
}
\end{tcolorbox}
\end{minipage}

\vspace{-2.5mm}
% --- KHỐI 3: SƠ ĐỒ TƯ DUY TIKZ ---
\noindent
\begin{tcolorbox}[
  enhanced,
  colback=orange!4!white,
  colframe=orange!80!black,
  arc=2mm, boxrule=0.8pt,
  left=2mm, right=2mm, top=1.5mm, bottom=1.5mm,
  title={\footnotesize\bfseries\color{white}\faSitemap\ SƠ ĐỒ TƯ DUY: MÔ HÌNH ĐOẠN THẲNG GẤP 3 LẦN}
]
\centering
\begin{tikzpicture}[
  seg/.style={draw=blue!80!black, fill=blue!15, rounded corners=1pt, minimum height=0.45cm, font=\scriptsize\bfseries}
]
  \node[anchor=east, font=\scriptsize\bfseries] at (0, 0.7) {Số bé ($4$):};
  \node[seg, minimum width=1.5cm] at (1.2, 0.7) {$4$};

  \node[anchor=east, font=\scriptsize\bfseries] at (0, 0) {Số lớn (gấp $3$ lần):};
  \node[seg, minimum width=1.5cm, fill=teal!20, draw=teal!80!black] at (1.2, 0) {$4$};
  \node[seg, minimum width=1.5cm, fill=teal!20, draw=teal!80!black] at (2.8, 0) {$4$};
  \node[seg, minimum width=1.5cm, fill=teal!20, draw=teal!80!black] at (4.4, 0) {$4$};

  \draw[thick, <->, purple!80!black] (0.45, -0.45) -- (5.15, -0.45) node[midway, below, font=\scriptsize\bfseries] {Số lớn = $4 \times 3 = \mathbf{12}$};
\end{tikzpicture}
\end{tcolorbox}

\vspace{-2.5mm}
% --- KHỐI 4: VÍ DỤ & MẸO TRÁNH BẪY ---
\noindent
\begin{minipage}[t]{0.49\textwidth}\vspace{0pt}
\begin{tcolorbox}[
  enhanced,
  colback=green!4!white,
  colframe=green!60!black,
  arc=2mm, boxrule=0.8pt,
  left=2mm, right=2mm, top=1.8mm, bottom=1.8mm,
  title={\footnotesize\bfseries\color{white}\faCheckCircle\ Ví dụ thực hành chuẩn mực}
]
{\footnotesize
\textbf{Bài toán:}\\
Năm nay An $7$ tuổi. Tuổi của mẹ gấp $5$ lần tuổi của An. Hỏi năm nay mẹ bao nhiêu tuổi?\\
\textbf{Bài giải:}\\
Năm nay tuổi của mẹ là:
$$7 \times 5 = \mathbf{35}\text{ (tuổi)}$$
\textbf{Đáp số:} $35$ tuổi.
}
\end{tcolorbox}
\end{minipage}
\hfill
\begin{minipage}[t]{0.49\textwidth}\vspace{0pt}
\begin{tcolorbox}[
  enhanced,
  colback=red!4!white,
  colframe=red!70!black,
  arc=2mm, boxrule=0.8pt,
  left=2mm, right=2mm, top=1.8mm, bottom=1.8mm,
  title={\footnotesize\bfseries\color{white}\faExclamationTriangle\ Mẹo tránh bẫy trắc nghiệm}
]
{\footnotesize
\begin{itemize}\setlength{\itemsep}{1.5pt}\setlength{\leftskip}{-2.5mm}
  \item \textbf{Bẫy nhầm từ khóa "Gấp" thành "Cộng":}\\
  Đề bài cho "gấp 4 lần", nhiều bạn đọc vội thành "nhiều hơn 4" nên làm phép cộng $7 + 4 = 11$ (SAI!). Phải làm phép nhân $7 \times 4 = 28$.
  \item \textbf{Bẫy thứ tự thừa số trong bài toán có lời văn:}\\
  Phải lấy số tuổi của 1 phần nhân với số lần: $7 \times 5$, không viết $5 \times 7$.
\end{itemize}
}
\end{tcolorbox}
\end{minipage}


\end{document}
